\section{Decision continuation spectra}
\label{sec:v20-spectrum}
The memory at an experimental cut has a different role from the memory
needed to reach that cut. We first identify its exact asymptotic value.
This will give both an obstruction and a construction for the same physical
audit, rather than a lower bound for an auxiliary binary problem.

\subsection{A charged training--validation interface}
Let $\Omega$ be a finite marked alphabet, $d=|\Omega|$, let
$B\subset\Delta(\Omega)$ be nonempty and compact, and fix
$Q\in\Delta(\Omega)$. A candidate $b\in B$ is fixed throughout the
experiment. Training consists of $n$ independent observations from $b$.
At the end of training the auditor retains a register $M\in[K]$.
Everything depending on the training observations or on an earlier random
seed must pass through this register. In particular, there is no external
copy of the training history and no uncharged persistent selector.

After this cut, fresh independent randomness may select an event $A$ using
$M$. The event is chosen before independent candidate and target validation
observations $X\sim b$ and $Y\sim Q$ are drawn. The score is
$\1_A(Y)-\1_A(X)$. The downstream register needed to keep $A$ and to
compute this score is distinct from the register at the specified cut.
The upstream and downstream widths will be stated separately in concrete
implementations. This is the reset interface of
Section~\ref{sec:v18-audit}, with one additional charged checkpoint.
The actual mark is part of a completed observation; it is not an input to
the private candidate's online transition rows.

Write $\mathcal H=[0,1]^\Omega$ and
\[
 f_h(b)=\sum_{w\in\Omega}(Q(w)-b(w))h(w),\qquad h\in\mathcal H.
\]
Every randomized event has a response $h(w)=\Pr(w\in A)$. Conversely,
every $h\in\mathcal H$ is the response of a random event: include each
atom independently with its specified probability. This is a finite
probability distribution on $2^\Omega$. Only its mean response matters
for expected signed score. The event itself must still be retained during
validation.

\begin{definition}[Decision continuation spectrum]
\label{def:v20-spectrum}
The $K$-state decision continuation value is
\begin{equation}\label{eq:v20-spectrum}
 \mathfrak V_K(B,Q)=
 \max_{h_1,\ldots,h_K\in\mathcal H}
 \min_{b\in B}\max_{m\le K} f_{h_m}(b).
\end{equation}
Let $v_{n,K}(B,Q)$ be the supremum of the worst-case expected score of
$n$-sample auditors at the preceding interface. Only the designated cut
has width $K$ in this definition; all other widths are finite but
unrestricted. For $c\in\R$, let $\kappa_c(B,Q)$ be the least $K$ for
which $\mathfrak V_K(B,Q)>c$, with value $+\infty$ when no such $K$
exists.
\end{definition}
The definition uses only the marked behavior image and its target, not a
parameterization of $B$ or a labeling of candidate states. No convex hull
of $B$ is taken. A response list is a list of conditional validation
continuations; it is not a mixture of feasible private candidates.

\begin{theorem}[Operational classification at a decision cut]
\label{thm:v20-spectrum}
All extrema in \eqref{eq:v20-spectrum} are attained. For every $n\ge1$,
\begin{equation}\label{eq:v20-rate}
 0\le \mathfrak V_K(B,Q)-v_{n,K}(B,Q)
 \le \min\left\{\sqrt{\frac{d-1}{n}},
                    \sqrt{\frac{2\log(2K)}{n}}\right\}.
\end{equation}
Consequently $v_{n,K}\to\mathfrak V_K$. Every auditor with $K$ states
at the designated cut, at any finite training length, has worst-case
score at most $\mathfrak V_K$. If $\mathfrak V_K>c$, some finite-sample
auditor with $K$ states at that cut has worst-case score greater than $c$.
Thus $\kappa_c$ is exactly the minimum cut width permitting such a strict
certificate with sufficiently many training samples.

Equivalently, $\kappa_c$ is the least number of response half-spaces
\begin{equation}\label{eq:v20-halfspaces}
 \{b\in B:f_h(b)>c\},\qquad h\in\mathcal H,
\end{equation}
that cover $B$. If two parameterized behavior families satisfy, uniformly
in their common parameter $q$,
\[
 \|b(q)-\widetilde b(q)\|_{\rm TV}\le\alpha,
 \qquad \|Q-\widetilde Q\|_{\rm TV}\le\beta,
\]
then
\begin{equation}\label{eq:v20-spectrumstable}
 |\mathfrak V_K(B,Q)-\mathfrak V_K(\widetilde B,\widetilde Q)|
 \le\alpha+\beta.
\end{equation}
\end{theorem}
\begin{proof}
For a fixed response list, the inner objective is continuous in $b$.
Its variation when the list changes is bounded uniformly over $b$ by
$2\max_m\|h_m-\widetilde h_m\|_\infty$. Its minimum over compact $B$
is therefore continuous in the list. Compactness of $\mathcal H^K$
proves attainment.

Condition on the retained state of any admissible auditor. Its conditional
post-cut event law depends only on that state and fresh randomness, and
has a response $h_m$. Independence of validation from training gives
\[
 g(b)=\sum_{m=1}^K \Pr_b(M=m)f_{h_m}(b)
                         \le\max_m f_{h_m}(b).
\]
Taking the minimum over $b$, and then the supremum over auditors, proves
the upper bound. This argument permits arbitrary stochastic training and
arbitrary finite upstream memory. An extra training-dependent decoder
would invalidate its first conditioning step and is excluded by the
interface, not silently supplied.

For the reverse bound fix an optimal list. Let $\widehat b$ be the
empirical measure and choose the first index maximizing
$f_{h_m}(\widehat b)$. Retain that index and discard the empirical table.
If $m_*$ maximizes the true score, then pointwise
\begin{equation}\label{eq:v20-regret}
 \max_m f_{h_m}(b)-f_{h_{\widehat m}}(b)
 \le2\max_{m\le K}|(\widehat b-b)h_m|.
\end{equation}
The maximum is at most $\|\widehat b-b\|_{\rm TV}$, since
$0\le h_m\le1$. Cauchy--Schwarz and the multinomial variances give
\[
 \E\|\widehat b-b\|_{\rm TV}^2
 \le\frac{d}{4n}\left(1-\sum_w b(w)^2\right)
 \le\frac{d-1}{4n}.
\]
This proves the first rate.

For completeness, if $V$ takes values in an interval of length one, then
$\log\E e^{s(V-\E V)}\le s^2/8$. Indeed the second derivative of the
logarithmic moment generating function is the variance under the
exponentially tilted law, at most $1/4$; its value and first derivative
vanish at zero. Integrating twice proves the inequality for both signs
of $s$. Independence consequently gives
$\E e^{s(\widehat b-b)h_m}\le e^{s^2/(8n)}$. For $s>0$,
\[
 \E\max_m|(\widehat b-b)h_m|
 \le\frac{\log(2K)}s+\frac{s}{8n}.
\]
Minimizing in $s$ gives $\sqrt{\log(2K)/(2n)}$ and proves the second
rate in \eqref{eq:v20-rate}. The empirical table has finitely many
values; the construction is therefore a finite-state auditor, with only
the chosen index surviving the specified cut. Real response coefficients
are allowed here; the physical construction below has deterministic
rational transitions.

A finite cover in \eqref{eq:v20-halfspaces} says that the continuous
function $b\mapsto\max_m f_{h_m}(b)$ is greater than $c$ everywhere on
compact $B$. Its attained minimum is then greater than $c$. This proves
the covering characterization in both directions. It also proves the
strict-threshold assertion; no claim about finite-sample attainment at
the boundary $c=\mathfrak V_K$ is needed.

Finally $|(Q-b)h-(\widetilde Q-\widetilde b)h|\le\alpha+\beta$
for every response. Taking maxima, minima over the common parameter set,
and maxima over lists preserves this uniform bound, proving
\eqref{eq:v20-spectrumstable}.
\end{proof}

\begin{corollary}[Compression by a finite decision dictionary]
\label{cor:v20-dictionary}
Suppose responses $h_1,\ldots,h_K$ satisfy
$\min_{b\in B}\max_m f_{h_m}(b)\ge c+\eta$ for $\eta>0$.
An empirical selector with $n>2\log(2K)/\eta^2$ achieves score greater
than $c$ and retains only $K$ states at the decision cut. If these
responses depend on a finite statistic $T:\Omega\to\Xi$, the training
selector can count $T(X)$ instead of $X$. A sufficient upstream register
has at most $\binom{n+|\Xi|-1}{|\Xi|-1}$ states at the last training
checkpoint, and the corresponding time-dependent count registers suffice
at earlier checkpoints.
\end{corollary}
\begin{proof}
Use the second bound in \eqref{eq:v20-rate}. Empirical expectations of
$h_m$ depend only on the histogram of $T(X)$. At time $t$, that histogram
has $\binom{t+|\Xi|-1}{|\Xi|-1}$ possibilities. Its update and the final
maximization are deterministic. No statistical sufficiency for estimating
all of $b$ is required; sufficiency for the specified responses is enough.
\end{proof}

This corollary supplies an intermediate principle between a full marked
histogram and a problem-specific compressed counter. The invariant is the
number and geometry of validation continuations needed to cover the
original candidate image at the desired score. It is not the algebraic
dimension of a chosen parameterization. The empirical upper bound need
not be sample-optimal, and a minimum decision-cut width is not a minimum
of the largest width over the entire acquisition schedule.

\subsection{Several feedback rows}
The classification does not require that the declared experiment have
only one feedback row. For the row-indexed interface of
Section~\ref{sec:v18-audit}, let
$\mathcal I=\{(j,A):A\subseteq\Omega_j\}$ and define
\[
 f_\lambda(b)=\sum_{i\in\mathcal I}\lambda_i f_i(b),
 \qquad \lambda\in\Delta(\mathcal I).
\]
A post-cut continuation is now a probability distribution on row--event
pairs; it is not an arbitrary test on a joint product of validation rows.
Put
\[
 \mathfrak V^{\rm rows}_K(B,Q)=
 \max_{\lambda_1,\ldots,\lambda_K\in\Delta(\mathcal I)}
       \min_{b\in B}\max_m f_{\lambda_m}(b).
\]
\begin{corollary}[Row-indexed classification]
\label{cor:v20-rows}
With $n$ training samples per row, the row-indexed decision-cut value
converges to $\mathfrak V^{\rm rows}_K$, and its deficit is at most
\[
 \min\{\sqrt{D/n},\sqrt{2\log(2K)/n}\},
 \qquad D=\sum_j(|\Omega_j|-1).
\]
The strict half-space cover characterization holds with $f_\lambda$.
The upper bound by $\mathfrak V^{\rm rows}_K$ applies also to any other
finite training schedule, including adaptive row selection.
\end{corollary}
\begin{proof}
Conditioning on the cut state gives the same convex-mixture upper bound
as before; its proof places no restriction on how that state was reached.
For empirical selection among an optimal list,
$|(\widehat b-b)\lambda_m|\le\max_j
\|\widehat b_j-b_j\|_{\rm TV}$. Sum the rowwise second moments to
obtain the first rate. To obtain the second, align the $r$th independent
training sample from each row into a tuple and write its response as
$\sum_{j,A}\lambda_{m,j,A}\1_A(X_{j,r})$. This random variable belongs
to $[0,1]$, and its sample mean is the empirical candidate mean for that
continuation. The logarithmic moment-generating argument in
Theorem~\ref{thm:v20-spectrum} therefore gives the same finite-list rate.
Each chosen row--event distribution can be drawn using fresh randomness
after the cut. Compactness and the covering argument are unchanged.
\end{proof}
